\documentclass[11pt]{article}

% --- Required style ---
\usepackage{series}

% --- Math packages ---
\usepackage{amsmath,amssymb,amsthm}
\usepackage{geometry}
\usepackage{hyperref}
\usepackage{booktabs}

\geometry{margin=1in}


% --- Metadata ---
\title{Theta Closure in Modal Triplet Theory I:
Gauge-Profile Targets from Multi-Loop Common-Scheme Transport}
\author{Peter Nero}
\date{July 2026}

\begin{document}
\maketitle
\begin{abstract}
We present a conditional $\Theta$--closure realization relating Standard Model
gauge couplings to internal geometric overlap integrals.  Measured source data
are transported with SMDR~v1.3 into the full-Standard-Model
$\overline{\mathrm{MS}}$ scheme at $Q=M_t$.  In the normalization used here,
the resulting profile targets are
\[
 I_2/I_1=0.5110273\pm0.0001231,\qquad
 I_3/I_1=0.158335\pm0.001098.
\]
They are calibrated profile coordinates, not first-principles predictions of
the gauge couplings.  We formulate the admissibility and spectral-gap
conditions under which an internal realization may reproduce these targets.
To address the normalization of nonabelian harmonics, we introduce a
high--coherence (twistor) corner with declared period and gauge-kinetic
normalizations. An $O(\lambda_Q^{-1})$ overlap estimate is conditional on
explicit projector and representative perturbation bounds.
The old $4.2$--$5~\mathrm{TeV}$ crossing and its identification with a physical
coherence, gap, or quantum-gravity scale are withdrawn.  The scale $Q=M_t$ is a
renormalization and matching convention only.  Explicit geometric realization
must therefore be re-executed against the new target pair.
\end{abstract}

\section*{Revision note for this edition}
\addcontentsline{toc}{section}{Revision note for this edition}
\begin{description}
\item[Supersedes.] \emph{Theta Closure in Modal Triplet Theory I: Gauge
Couplings from Internal Geometry}, first edition.
\item[Reason.] The former target used an obsolete one-loop few-TeV crossing,
treated calibrated ratios as predictions, and combined unit-$L^2$
normalization with overlap definitions that would force trivial unit values.
\item[Resolution.] Version~2 uses weighted gauge-kinetic coefficients and
SMDR~v1.3 common-scheme targets at $Q=M_t$, with explicit provenance and
conditional projector-error assumptions.
\item[Retained result.] Gauge-coupling ratios can still be posed as geometric
overlap targets in a normalized high-coherence realization.
\item[Remaining boundary.] MTT must independently select and execute the
internal geometry before these calibrated targets become predictions.
\end{description}

\section{Purpose and scope}

This paper presents a conditional map from a common-scheme Standard Model gauge
profile to internal overlap targets in Modal Triplet Theory (MTT).

The goals are deliberately narrow:
\begin{enumerate}
\item Derive gauge--coupling expressions from internal overlap integrals.
\item Extract numerical $\Theta$--profile targets from measured couplings.
\item Test whether a nonempty admissible $\Theta$--region satisfies the MTT gap conditions.
\item State all assumptions and calibration choices explicitly.
\end{enumerate}

No claim of uniqueness, zero-knob gauge-coupling prediction, or full
unification is made.  The coupling data enter upstream of the profile targets;
agreement of a fitted geometry with those same targets is a realization or
round-trip test, not a held-out prediction.

\section{Standing assumptions of Modal Triplet Theory}

We collect here the assumptions used throughout.
They are standard within the MTT corpus and are not re-derived.

\begin{assumption}[Bounded internal geometry]
The internal bundles $B_n$ possess bounded geometry and admit commuting vertical Laplacians.
\end{assumption}

\begin{assumption}[Uniform spectral gap]
There exists a strictly positive lower bound
\[
\lambda_{\ast} > 0
\]
separating coherent from noncoherent internal modes.
\end{assumption}

\begin{assumption}[Bounded coherent projection]
The joint Riesz projector
\[
\Pi = \Pi_{B_1}\Pi_{B_2}\Pi_{B_3}
\]
is bounded on $H^1$.
\end{assumption}

\begin{assumption}[Fundamental Contractivity Condition (FCC)]
The projected flow $F=\Pi\circ\Phi_\tau$ is contractive on the coherent sector.
\end{assumption}

\begin{remark}
If these assumptions fail, this paper loses controlled access to the stated
coherent effective encoding.  No claim is made that the underlying
mathematical configuration, probabilities, records, or observers cease to
exist in every description.
\end{remark}

\subsection{Spectral gap as a controlled-description gate}

A uniform spectral gap is assumed throughout this work as a sufficient
condition for a stable finite coherent sector. If the selected cluster loses
separation, the corresponding Riesz projector and truncation estimates may
lose regularity, so this paper no longer controls the stated effective
encoding. This does not imply that amplitudes, probabilities, records, or an
underlying mathematical description cease to exist in every formulation.

The present work does not derive the gap dynamically or prove that every
physical theory must possess this particular finite-cluster description. It
tests consequences conditional on the declared gap and projector hypotheses.

\section{Internal operator structure}

The physical dimensional-reduction formula is posed on one compact
six-manifold $X_6$ with three compatible vertical sectors. A single Hermitian
line bundle or principal $U(1)$ bundle $L_{\mathrm{shared}}\to X_6$ supplies
the common phase/holonomy datum. It acts in every sector and is counted once.

For the calibrated Route--A computation only, write $B_n$ for the effective
support and measure used by sector $n$. Locally these supports may be modeled
using:
\begin{itemize}
\item a shared-circle phase normalization for the abelian lane,
\item a round two-dimensional lens-base representative for the weak lane,
\item a compact Heisenberg nilmanifold representative for the color lane.
\end{itemize}
The three $B_n$ are not asserted to be simultaneous Cartesian factors of
$X_6$, and no literal
$S^1_{\mathrm{shared}}\times L(3,1)\times\mathrm{Nil}_3$ compactification is
used. The Lens--Nil data form an auxiliary calibrated operator model; the
current global physical candidate is the distinct q79/Fu--Yau branch.

\section{Gauge couplings from internal overlaps}

\subsection{Definition of gauge-kinetic overlaps}

\begin{definition}[Sector gauge-kinetic coefficient]
For each gauge factor $a=1,2,3$, fix a sector representative $\omega_a$ by a
declared pointwise, period, or cohomological normalization and define
\[
 I_a := \int_{B_a}
 w_a\,\langle \omega_a,\omega_a\rangle\,d\mu_{B_a}.
\]
Here $w_a>0$ contains the induced gauge-kinetic density, including any
integrated spectator directions and trace factors. The representative is not
simultaneously normalized to unit $L^2(B_a,w_a\,d\mu)$; doing so would make
$I_a=1$ by definition and erase the coupling information.
\end{definition}

\subsection{Reduction of the Yang--Mills action}

\begin{proposition}[Conditional coupling--overlap relation]
Assume a ten-dimensional Yang--Mills sector with common coefficient
$g_{10}$, a consistent truncation to mutually orthogonal four-dimensional
gauge zero modes, and the gauge-kinetic weights and normalizations in the
preceding definition. Then
\[
\frac{1}{g_a^2} = \frac{1}{g_{10}^2}\, I_a.
\]
\end{proposition}

\begin{proof}
Start from the ten-dimensional Yang--Mills action
\[
S_{10} = \frac{1}{2g_{10}^2}\int \mathrm{Tr}(F_{AB}F^{AB})\, d\mathrm{vol}_{10}.
\]

Expand the gauge field along the declared internal sector modes
\[
A_\mu(x,u) = A_\mu^{(a)}(x)\,\omega_a(u),
\]

Integrating the weighted internal gauge-kinetic density and using orthogonality
of the retained modes yields
\[
S_4 = \frac{1}{4g_a^2}\int F_{\mu\nu}^{(a)}F^{(a)\mu\nu}\, d\mathrm{vol}_4,
\]
with
\[
\frac{1}{g_a^2}
=\frac{1}{g_{10}^2}
\int_{B_a}w_a\langle\omega_a,\omega_a\rangle d\mu_{B_a}.
\]
\end{proof}

\subsection{Normalization conventions}

We adopt the following conventions:

\begin{enumerate}
\item \textbf{Generator normalization:}
\[
\mathrm{Tr}(T^a T^b)=\frac{1}{2}\delta^{ab}.
\]

\item \textbf{Hypercharge normalization (GUT-normalized):}
\[
g_1=\sqrt{\frac{5}{3}}\, g',
\qquad
\alpha_1=\frac{5}{3}\alpha_Y.
\]
\end{enumerate}

With these conventions, the normalization constants satisfy
\[
N_1=N_2=N_3=1,
\]
and all group-theoretic factors are absorbed into the definition of $g_1$.

\section{Summary of this section}

At this point we have established:

\begin{itemize}
\item Gauge couplings descend from weighted internal gauge-kinetic coefficients
under the stated higher-dimensional action and truncation assumptions.
\item The relation $1/g_a^2=I_a/g_{10}^2$ is exact within that declared
reduction ansatz.
\item Ratios of couplings are ratios of those coefficients only when the same
$g_{10}$ and compatible generator conventions apply to all three sectors.
\end{itemize}

In the next part we will:
\begin{enumerate}
\item Compute $g_1,g_2,g_3$ explicitly from experimental inputs.
\item Run them to a common matching scale.
\item Extract numerical $\Theta$–constraints.
\end{enumerate}
% ===========================
\section{Experimental inputs and gauge couplings at $M_Z$}
% ===========================

In this section we compute the Standard Model gauge couplings explicitly
from experimental inputs at the $Z$-pole and prepare them for comparison
with the internal overlap formulation.

\subsection{Electroweak inputs}

We use the following $\overline{\mathrm{MS}}$-scheme inputs at $\mu=M_Z$:
\begin{equation}
\alpha(M_Z)^{-1} = 127.95,
\qquad
\sin^2\theta_W(M_Z) = 0.23122,
\qquad
\alpha_s(M_Z) = 0.1179.
\label{eq:EWinputs}
\end{equation}

These values are representative of current PDG averages and are sufficient
for the present analysis.

\subsection{Definition of gauge couplings}

Define the electromagnetic coupling
\begin{equation}
e := \sqrt{4\pi \alpha(M_Z)}.
\end{equation}

Then the electroweak couplings are
\begin{equation}
g_2 = \frac{e}{\sin\theta_W},
\qquad
g' = \frac{e}{\cos\theta_W},
\end{equation}
and the GUT-normalized hypercharge coupling is
\begin{equation}
g_1 = \sqrt{\frac{5}{3}}\, g'.
\end{equation}

The strong coupling is
\begin{equation}
g_3 = \sqrt{4\pi \alpha_s(M_Z)}.
\end{equation}

\subsection{Explicit numerical evaluation}

From \eqref{eq:EWinputs}:
\[
\alpha(M_Z) \approx 0.007816,
\qquad
e = \sqrt{4\pi\alpha} \approx 0.3135.
\]

Also,
\[
\sin\theta_W = \sqrt{0.23122} \approx 0.4809,
\qquad
\cos\theta_W \approx 0.8768.
\]

Thus,
\[
g_2(M_Z) = \frac{0.3135}{0.4809} \approx 0.6517,
\]
\[
g'(M_Z) = \frac{0.3135}{0.8768} \approx 0.3576,
\]
\[
g_1(M_Z) = \sqrt{\frac{5}{3}}\times 0.3576 \approx 0.4614.
\]

For QCD,
\[
g_3(M_Z) = \sqrt{4\pi\times 0.1179} \approx 1.2172.
\]

\begin{table}[h]
\centering
\begin{tabular}{lccc}
\toprule
Scale & $g_1$ (GUT) & $g_2$ & $g_3$ \\
\midrule
$M_Z$ & $0.4614$ & $0.6517$ & $1.2172$ \\
\bottomrule
\end{tabular}
\caption{Gauge couplings at $\mu=M_Z$ computed from \eqref{eq:EWinputs}.}
\end{table}

% ===========================

% ===========================
\section{Selected multi-loop common-scheme transport}
% ===========================

The former one-loop evolution to $5~\mathrm{TeV}$ is not used in this
revision.  In particular, the Standard Model equations do not place the
$g_1=g_2$ crossing there.  We instead use the selected SMDR~v1.3 transport,
which maps fifteen measured source coordinates to the full non-decoupled
Standard Model in the tadpole-free pure $\overline{\mathrm{MS}}$ scheme at
\[
 Q=M_t=172.5590883453979~\mathrm{GeV}.
\]
This is a scheme and comparison scale, not an MTT coherence scale.

The accepted gauge rows are
\begin{align}
 g_Y(Q)&=0.3585945042\pm0.0000307251,\\
 g_2(Q)&=0.6475986708\pm0.0000287665,\\
 g_3(Q)&=1.163427409\pm0.004036156.
\end{align}
With the paper's GUT normalization,
\[
 g_1(Q)=\sqrt{\frac53}\,g_Y(Q)
       =0.4629435143\pm0.0000396660.
\]
These rows are outputs of a common multi-loop matching/running map.  Their
covariance is propagated from the declared diagonal measured-input profile;
an official joint likelihood spanning all fifteen source coordinates is not
claimed.

% ===========================
\section{Numerical $\Theta$--profile targets from gauge data}
% ===========================

Under the coupling--overlap relation already proved conditionally above,
\[
 \frac{I_b}{I_a}=\left(\frac{g_a}{g_b}\right)^2.
\]
Consequently the selected common-scheme profile gives
\begin{equation}
\boxed{
 \frac{I_2}{I_1}=0.5110273\pm0.0001231,
 \qquad
 \frac{I_3}{I_1}=0.158335\pm0.001098.
}
\label{eq:targets}
\end{equation}
The propagated covariance of the two ratios is
$-6.1892\times10^{-9}$, corresponding to correlation $-0.04578$.

These are experimentally anchored profile targets.  A geometry adjusted to
reproduce them is a calibrated realization.  A held-out prediction would
require fixing that geometry without these gauge rows and then computing an
observable not used in source, branch, scale, or model selection.

\section{Transition to the geometric $\Theta$-problem}

The next step is to determine whether the internal geometry of MTT
admits overlap integrals $I_1,I_2,I_3$ satisfying the above ratios
while maintaining the baseline spectral gap $\lambda_{\ast}=0.25$.

This problem is addressed in the next section.
% ===========================
\section{Baseline internal geometry and spectral gap bounds}
% ===========================

We now specify the auxiliary operator models used for the Route--A gap check
and collect their assumed spectral inequalities.

\subsection{Auxiliary metric supports}

The three vertical operators are tested on effective supports $\Sigma_n$:
\begin{itemize}
\item $\Sigma_1$ is an auxiliary circle,
\item $\Sigma_2$ is a two-dimensional lens-base model,
\item $\Sigma_3$ is a compact nilmanifold model.
\end{itemize}
The common $U(1)$ phase/holonomy line acts across all three supports and is
counted once. The notation does not assert three products
$S^1_{\mathrm{cen}}\times\Sigma_n$ or identify these supports with coordinate
factors of the selected q79 compactification.

The auxiliary metric ansatz has parameters:
\begin{itemize}
\item $R_1$ -- radius of $\Sigma_1$,
\item $R_{\mathrm{lens}}$ -- effective lens-base radius,
\item $f_2$ -- lens-base scale factor,
\item $h_0$ -- Cheeger lower-bound parameter for the nil model.
\end{itemize}

\subsection{Spectral gap bounds}

The following are model-specific lower bounds on the smallest nonzero
eigenvalues of the auxiliary vertical Laplacians.

\paragraph{Circle factor $\Sigma_1$.}
For a circle of radius $R_1$,
\begin{equation}
\lambda_{\Sigma_1} \sim \frac{1}{R_1^2}.
\label{eq:gapSigma1}
\end{equation}

\paragraph{Lens factor $\Sigma_2$.}
For the lens sheet with effective radius $f_2R_{\mathrm{lens}}$,
\begin{equation}
\lambda_{\mathrm{lens}} \ge \frac{2}{(f_2R_{\mathrm{lens}})^2}.
\label{eq:gapLens}
\end{equation}

\paragraph{Nil factor $\Sigma_3$.}
For the nil leaf, the spectral gap is controlled by the Cheeger constant:
\begin{equation}
h \ge h_0 > 0
\quad\Rightarrow\quad
\lambda_{\mathrm{nil}} \ge \frac{h_0^2}{4}.
\label{eq:gapNil}
\end{equation}

\subsection{Baseline numerical values}

In the worked numerical baseline of the Foundation,
\begin{equation}
h_0 = 1
\quad\Rightarrow\quad
\lambda_{\mathrm{nil}} = \frac{1}{4} = 0.25.
\end{equation}

We enforce the \emph{baseline lower bound} $\lambda_{\ast}\ge 0.25$ using $h_0=1$.
This sets a conservative admissibility floor; explicit realizations of the nil sector may
exceed this bound and need not saturate it.


% ===========================
\section{Geometric overlap model}
% ===========================

We now introduce a minimal, computable auxiliary coefficient model. It is not
obtained by multiplying three sector supports or by appending the shared circle
to each support. Instead, all spectator integrations, mode normalizations, and
gauge-kinetic weights are absorbed into three positive coefficients.

\subsection{Auxiliary coefficient ansatz}

At the Route--A calibration tier, parameterize the coefficients as
\begin{align}
I_1 &= 2\pi R_1,
\label{eq:I1}\\
I_2 &= \kappa_\ell (f_2R_{\mathrm{lens}})^2,
\label{eq:I2}\\
I_3 &= \kappa_n\,\mathcal{S}_n,
\label{eq:I3}
\end{align}
This is an ansatz for the weighted coefficients defined above, not a theorem
equating the norm of an $L^2$-normalized mode with
a geometric volume. A selected compactification must emit the corresponding
weights and representatives on one $X_6$ before these coefficients become
source-derived.
where:
\begin{itemize}
\item $\kappa_\ell$ is the unit-area normalization of the lens sheet,
\item $\kappa_n\,\mathcal{S}_n$ is the effective nil overlap weight,
\item $\mathcal{S}_n$ is not assumed equal to a simple metric area.
\end{itemize}

This separation is required because the nil spectral gap
\eqref{eq:gapNil} is controlled by $h_0$, not directly by metric scale.

% ===========================
\section{Solving the $\Theta$--system}
% ===========================

We now solve the overlap constraints from Section~4 of Part~II
together with the spectral gap bounds.

\subsection{Overlap constraints}

From \eqref{eq:targets} and \eqref{eq:I1}--\eqref{eq:I3}:
\begin{align}
\frac{I_2}{I_1}
&=
\frac{\kappa_\ell (f_2R_{\mathrm{lens}})^2}{2\pi R_1}
= 0.5110273,
\label{eq:Theta1}\\
\frac{I_3}{I_1}
&=
\frac{\kappa_n\mathcal{S}_n}{2\pi R_1}
= 0.158335.
\label{eq:Theta2}
\end{align}

Solving \eqref{eq:Theta1} gives
\begin{equation}
(f_2R_{\mathrm{lens}})^2
=
\frac{0.5110273\cdot 2\pi}{\kappa_\ell}\,R_1
=
\frac{3.210879}{\kappa_\ell}\,R_1,
\label{eq:RlSolve}
\end{equation}
hence
\begin{equation}
R_{\mathrm{lens}}
=
\frac{\sqrt{(3.210879/\kappa_\ell)\,R_1}}{f_2}.
\label{eq:RlFinal}
\end{equation}

From \eqref{eq:Theta2}:
\begin{equation}
\kappa_n\mathcal{S}_n
=
0.158335\cdot 2\pi\,R_1
=
0.994849\,R_1.
\label{eq:SnSolve}
\end{equation}

\subsection{Imposing gap inequalities}

We require:
\begin{align}
\lambda_{\Sigma_1} &\ge 0.25,
\label{eq:gap1}\\
\lambda_{\mathrm{lens}} &\ge 0.25,
\label{eq:gap2}\\
\lambda_{\mathrm{nil}} &\ge 0.25.
\label{eq:gap3}
\end{align}

\paragraph{Nil.}
Set $h_0=1$, so \eqref{eq:gapNil} gives $\lambda_{\mathrm{nil}}=0.25$.

\paragraph{Circle.}
From \eqref{eq:gapSigma1} and \eqref{eq:gap1},
\[
\frac{1}{R_1^2}\ge 0.25
\quad\Rightarrow\quad
R_1\le 2.
\]

\paragraph{Lens.}
From \eqref{eq:gapLens} and \eqref{eq:gap2},
\[
\frac{2}{(f_2R_{\mathrm{lens}})^2}\ge 0.25
\quad\Rightarrow\quad
(f_2R_{\mathrm{lens}})^2\le 8.
\]

Substitute \eqref{eq:RlSolve}:
\[
\frac{3.210879}{\kappa_\ell}\,R_1 \le 8
\quad\Rightarrow\quad
R_1 \le \frac{8\kappa_\ell}{3.210879}\approx 2.492\,\kappa_\ell.
\]

\subsection{Existence theorem}

\begin{theorem}[Existence of an auxiliary calibrated $\Theta$--region]
At the selected common-scheme profile point $Q=M_t$, choose
\[
h_0=1,\qquad \kappa_\ell=1,\qquad f_2\ge 1.
\]
Then for every $R_1\in(0,2]$ there exists a choice of
$R_{\mathrm{lens}}$ and $\mathcal{S}_n$ given by
\eqref{eq:RlFinal} and \eqref{eq:SnSolve} such that:
\begin{enumerate}
\item the gauge overlap constraints
$\frac{I_2}{I_1}=0.5110273$ and $\frac{I_3}{I_1}=0.158335$
are satisfied;
\item the spectral gap inequalities
$\lambda_{\Sigma_1},\lambda_{\mathrm{lens}},\lambda_{\mathrm{nil}}\ge 0.25$
hold;
\item the minimal spectral gap remains $\lambda_{\ast}=0.25$.
\end{enumerate}
\end{theorem}

\begin{proof}
Let $R_1\in(0,2]$. Then $\lambda_{\Sigma_1}\ge 0.25$.
With $\kappa_\ell=1$ and $f_2\ge 1$, \eqref{eq:RlFinal} implies
$(f_2R_{\mathrm{lens}})^2=3.210879R_1\le 6.421758<8$, hence
$\lambda_{\mathrm{lens}}>0.25$.
With $h_0=1$, $\lambda_{\mathrm{nil}}=0.25$.
The overlap constraints are satisfied by construction.
\end{proof}

\section{Interpretation}

The existence theorem shows only that the measured Standard Model gauge
profile can be represented inside this auxiliary coefficient ansatz without
violating its assumed dimensionless gap inequalities. The choice $h_0=1$
saturates the declared nil lower bound in this ansatz; it does not prove that
the selected q79 geometry has this value or that the nil lane physically
controls the onset of noncoherent behavior.

\begin{remark}
The quantities $\kappa_n$ and $\mathcal{S}_n$ encode nontrivial
nilmanifold structure and are not free once explicit harmonic
representatives are specified.
\end{remark}

\begin{remark}[Coherent sector stability]
All results in this work are conditional on the stability of the coherent
sector.
If the coherent sector were dynamically unstable, the theory would not merely
predict different numerical values; rather, it would fail to support persistent
observables, reproducible measurements, or probabilistic interpretation.
Accordingly, coherent sector stability is treated here as a prerequisite for
physics rather than a phenomenological assumption.
\end{remark}

% ===========================
\section{Massless gauge-sector normalization via the twistor corner}
% ===========================

The remaining technical ambiguity in the $\Theta$--closure analysis concerns
the explicit definition and normalization of the nonabelian gauge-kinetic
coefficients $I_2$ and $I_3$ associated with the $SU(2)$ and $SU(3)$ sectors.
In this section we state a well-typed conditional definition in the
high--coherence (twistor) corner of Modal Triplet Theory.

\begin{remark}[Twistor corner as a computational device]
In this paper the twistor corner is used solely as a \emph{computational normalization regime}
for the massless coherent gauge sector, providing canonical harmonic representatives and
gap-controlled error bounds. No claim is made that Nature resides exactly in this corner,
nor that twistor space is fundamental.
\end{remark}


\subsection{High--coherence regime and admissible truncation}

The twistor formulation of MTT establishes the existence of a high--coherence
regime in which:
\begin{enumerate}
\item the spectral gap $\lambda_Q$ separating coherent and noncoherent modes
is parametrically large;
\item the effective generator admits a Schur--Feshbach reduction onto a
finite--dimensional coherent subspace;
\item truncation errors are bounded by explicit operator--norm estimates
of order $O(\lambda_Q^{-1})$.
\end{enumerate}

\begin{assumption}[Twistor-corner admissibility]
We assume that the $\Theta$--matching scale $\mu_\Theta$ lies within the
high--coherence regime for the massless gauge sector, so that the
Schur--Feshbach reduction is valid with remainder bounded by
$O(\lambda_Q^{-1})$.
\end{assumption}

\begin{remark}
This assumption invokes no ultraviolet completion and no string-theoretic
structure; it relies only on the coherence and gap hypotheses already stated
in Sections~2 and~10.
\end{remark}

\subsection{Canonical definition of the massless gauge subspace}

Let $\mathcal H$ denote the full internal Hilbert space of modal fields.
Define $\mathcal H_0\subset\mathcal H$ as the coherent massless gauge subspace
selected by the twistor corner, characterized by:
\begin{enumerate}
\item self--duality of the Yang--Mills curvature;
\item vanishing effective mass under the reduced generator;
\item holomorphic encoding via the Ward correspondence.
\end{enumerate}

\begin{definition}[Twistor-corner massless sector]
$\mathcal H_0$ is defined as the image of the coherent projector $P_0$
associated with the self--dual Yang--Mills corner of the MTT configuration
space.
\end{definition}

This definition is canonical within the admissible slab and does not depend
on additional choices.

\subsection{Schur--Feshbach reduction and error control}

Let $L$ be the full internal generator and write
\[
L =
\begin{pmatrix}
P_0 L P_0 & P_0 L Q \\
Q L P_0 & Q L Q
\end{pmatrix},
\qquad Q=1-P_0.
\]

The effective generator on $\mathcal H_0$ is given by the Schur--Feshbach map
\[
L_{\mathrm{eff}}
=
P_0 L P_0
-
P_0 L Q (Q L Q)^{-1} Q L P_0.
\]

The twistor analysis proves the operator bound
\[
\|P_0 L Q (Q L Q)^{-1} Q L P_0\|
\;\le\;
C\,\lambda_Q^{-1},
\]
for some constant $C$ independent of infrared data.

\subsection{Harmonic representatives and normalization}

\begin{definition}[Twistor-corner harmonic representatives]
For $a=2,3$, fix a cohomology class together with its period normalization,
the internal metric, and a gauge condition. Let $\omega_a^{(0)}$ denote its
harmonic representative in $\mathcal H_0$. Hodge theory fixes the harmonic
representative of the fixed class; the period condition, rather than unit
$L^2$ normalization, fixes its scale.
\end{definition}

\begin{remark}
The abelian overlap $I_1$ remains fixed by the explicit $S^1_{\mathrm{cen}}$
normalization of Appendix~G and is not modified by the twistor construction.
\end{remark}

\subsection{Leading--order overlap integrals}

\begin{definition}[Leading--order overlaps]
The leading--order nonabelian overlap integrals are defined by
\[
I_a^{(0)} := \int_{B_a}
w_a^{(0)}
\langle \omega_a^{(0)}, \omega_a^{(0)} \rangle \, d\mu_{B_a},
\qquad a=2,3.
\]
\end{definition}

These depend on the massless coherent geometry, the fixed periods, and the
gauge-kinetic weights. The twistor representation alone does not select those
normalization data.

\subsection{Control of overlap deviations}

Let $I_a$ denote the full overlap integrals defined in Section~4.

\begin{lemma}[Conditional overlap stability under coherent truncation]
Assume the full and leading projectors, period-normalized representatives,
weights, and measures obey uniform perturbation estimates of order
$\lambda_Q^{-1}$ in norms controlling the displayed integrals. Then there
exist constants $C_a>0$ such that
\[
|I_a - I_a^{(0)}| \;\le\; C_a\,\lambda_Q^{-1},
\qquad a=2,3.
\]
\end{lemma}

\begin{proof}
Expand the difference of the two weighted quadratic integrals into the
representative, weight, and measure perturbations. Cauchy--Schwarz and the
assumed uniform bounds control each term by its corresponding
$O(\lambda_Q^{-1})$ estimate. Their finite sum gives the stated constant
$C_a$. The Schur--Feshbach bound alone would not supply all of these
coefficient estimates.
\end{proof}

\subsection{Twistor-corner $\Theta$--closure at leading order}

Combining the above with the experimentally extracted targets at the selected common-scheme point $Q=M_t$:
\[
\frac{I_2}{I_1} = 0.5110273,
\qquad
\frac{I_3}{I_1} = 0.158335,
\]
we obtain:

\begin{theorem}[Twistor-corner $\Theta$--closure condition]
If the twistor-corner admissibility assumption holds at $\mu_\Theta$, then
$\Theta$--closure at leading order reduces to the computable conditions
\[
\frac{I_2^{(0)}}{I_1} = 0.5110273,
\qquad
\frac{I_3^{(0)}}{I_1} = 0.158335,
\]
with controlled corrections of order $O(\lambda_Q^{-1})$.
\end{theorem}

\begin{remark}
Failure of these relations falsifies either the baseline internal geometry
or the twistor-corner admissibility hypothesis.
\end{remark}

\subsection{Scope}

This section does not claim that the twistor corner describes the full
physical gauge sector. It provides a conditional computational corner in which
$I_2$ and $I_3$ are well defined after the metric, periods, weights, and
perturbation bounds have been supplied.

% ===========================

% ===========================
\section{Scale separation and withdrawn legacy calibration}
% ===========================

The scale $Q=M_t$ used above is a renormalization/matching coordinate.  It does
not determine the physical internal gap, compactification radius, proper-time
filter, Planck scale, Hubble scale, or primordial tensor amplitude.  The former
identifications
\[
 Q\sim E_{\mathrm{gap,min}}\sim\tau_0^{-1/2}
\]
were additional calibrations tied to the invalid $4.2$--$5~\mathrm{TeV}$
crossing and are withdrawn.

The dimensionless spectral inequalities in the preceding sections remain
conditional geometric statements.  Converting them to SI or external energy
units requires an independently selected action normalization and a theorem
relating the internal operator spectrum to physical propagator poles or
response scales.  No cosmological or quantum-gravity bound is inferred in this
paper.


% ===========================
\section{Falsifiability and status}
% ===========================

At the profile tier, the direct test is whether a specified internal geometry,
with its measure and harmonic normalization fixed independently of the gauge
targets, emits the two ratios in \eqref{eq:targets} within the propagated
covariance.  If the geometry is fitted to those ratios, the result is instead
a realization and round-trip consistency check.

The old cosmological and Gaussian quantum-gravity discriminators are removed
because their physical scale identification is not derived here.  The current
broader MTT repository closes embedded renormalized-Standard-Model equivalence
at the adopted one-shared-physical-primitive/profile standard.  That successor
does not turn the present overlap targets into zero-knob predictions and does
not derive standard perturbative quantization from MTT.

\section{Conclusion}

The coupling--overlap formula survives as a conditional dimensional-reduction
identity.  Its numerical implementation is now placed in one selected
multi-loop common scheme at $Q=M_t$, yielding
\[
 I_2/I_1=0.5110273\pm0.0001231,\qquad
 I_3/I_1=0.158335\pm0.001098.
\]
The full covariance provenance is retained, and the absence of a public joint
fifteen-coordinate likelihood is stated explicitly.

The former $5~\mathrm{TeV}$ crossing, the geometry calibrated to its ratios,
and the identification of that scale with an internal gap, proper-time cutoff,
or cosmological scale are not results of this revision.  Papers II--V must be
re-executed or reclassified accordingly.  This paper therefore supplies a
reproducible profile target and a precise geometric test, not a unique
first-principles prediction of gauge couplings or internal geometry.

\section*{References}
% ===========================

\begin{enumerate}
\item Particle Data Group (PDG), Review of Particle Physics (for $M_Z$, $\alpha(M_Z)$, $\sin^2\theta_W(M_Z)$, $\alpha_s(M_Z)$).
\item M. E. Machacek and M. T. Vaughn, ``Two-loop renormalization group equations in a general quantum field theory,'' Nucl. Phys. B222 (1983) 83--103.  % for RGE conventions
\item T. Kato, \emph{Perturbation Theory for Linear Operators}, Springer (1995). % for spectral projection continuity
\item Peter Nero, \emph{Modal Triplet Theory: Foundation} (MTT corpus).
\item Peter Nero, \emph{Modal Triplet Theory: Quantum Amplitudes from Modal Geometry} (MTT corpus).
\item Peter Nero, \emph{Twistor Encodings as High-Coherence Limits of Modal Triplet Theory} (MTT corpus).
\end{enumerate}

\end{document}
